% Part: computability
% Chapter: computability-theory
% Section: ce-closed-cup-cap

\documentclass[../../../include/open-logic-section]{subfiles}

\begin{document}

\olfileid[fr]{cmp}{thy}{clo}
\olsection[Réunion et intersection d'ensembles effectivement énumérables]{Stabilité des ensembles effectivement énumérables par réunion et intersection}

Le théorème suivant donne deux propriétés de
stabilité des ensembles effectivement énumérables.

\begin{thm}
Si $A$ et~$B$ sont effectivement énumérables,
$A\cap B$ et $A\cup B$ le sont aussi.
\end{thm}

\begin{proof}
Le \olref[eqc]{thm:ce-equiv} permet d'utiliser
diverses caractérisations de ces ensembles.
Nous allons donner plusieurs démonstrations.
Si $A$ ou~$B$ est vide, l'intersection est vide
et la réunion est l'autre ensemble : le résultat
est immédiat. Supposons donc $A$ et~$B$ non vides
pour les constructions qui suivent.\footnote{Complément éditorial : les deux premières constructions source utilisent des fonctions totales qui énumèrent $A$ et $B$, ce qui exige qu'ils soient non vides. Nous séparons les cas vides avant de les appliquer.}

Pour la première preuve, supposons que $f$ énumère
$A$ et que $g$ énumère~$B$, avec $f$ et~$g$
calculables. Posons
\begin{align*}
h(x) & = \umin{y}{(f(y) = x \lor g(y) = x)} \text{ et}\\
j(x) & = \umin{y}{(f((y)_0) = x \land g((y)_1) = x)}.
\end{align*}
Alors $A\cup B$ est le domaine de~$h$ et
$A\cap B$ celui de~$j$.

\begin{explain}
En termes de calcul, si des procédures énumèrent
$A$ et~$B$, on peut semi-décider l'appartenance
d'un élément~$x$ à $A\cup B$ en le cherchant dans
l'une ou l'autre énumération; pour $A\cap B$, on
le cherche simultanément dans les deux.
\end{explain}

Pour la deuxième preuve, supposons encore que
$f$ énumère $A$ et que $g$ énumère~$B$. Posons
\[
k(x) = \begin{cases}
f(x/2) & \text{si $x$ est pair} \\
g((x-1)/2) & \text{si $x$ est impair.}
\end{cases}
\]
La fonction $k$ énumère $A\cup B$ en alternant
les valeurs fournies par $f$ et~$g$. Pour
$A\cap B$, la construction est plus délicate.
Si l'intersection est vide, elle est effectivement
énumérable par définition. Sinon, choisissons
$c\in A\cap B$ et définissons $l$ par
\[
l(x) = \begin{cases}
f((x)_0) & \text{si $f((x)_0) = g((x)_1)$} \\
c & \text{sinon.}
\end{cases}
\]
La fonction $l$ parcourt les paires de valeurs
produites par $f$ et~$g$, et renvoie chaque valeur
commune; en l'absence de coïncidence, elle renvoie~$c$.

Pour la dernière preuve, supposons que $A$ soit
le \emph{domaine} d'une fonction calculable partielle
$m(x)$ et $B$ celui d'une fonction calculable
partielle $n(x)$. Alors $A\cap B$ est le domaine
de la fonction partielle $m(x)+n(x)$.

\begin{explain}
En termes de calcul, si $A$ est l'ensemble des
entrées sur lesquelles $m$ s'arrête et $B$ celui
des entrées sur lesquelles $n$ s'arrête, leur
intersection est l'ensemble des entrées où les
deux procédures s'arrêtent.
\end{explain}

Exprimer $A\cup B$ comme domaine d'une fonction
calculable partielle exige de simuler $m$ et~$n$
en parallèle. Soient $d$ et~$e$ des indices de
$m$ et~$n$, respectivement, donc
$m\simeq\cfind{d}$ et $n\simeq\cfind{e}$.
Alors $A\cup B$ est le domaine de la fonction
\[
p(x) = \umin{y}{(T(d,x,y) \lor T(e,x,y))}.
\]
\begin{explain}
Sur l'entrée~$x$, la fonction $p$ cherche un
calcul qui s'arrête pour $m$ ou pour~$n$, et
est définie si elle en trouve un.
\end{explain}
\end{proof}

\end{document}
